%!TEX root = ../main.tex
\section{Proof of Theorem \ref{thm:lqr-variance-T-compose}}
\label{sec:app-proof-lqr-variance-T-compose}

\begin{proof}
Recall the policy gradient estimators
\[
\widehat G_K
=\sum_{t=0}^{T-1}\Sigma^{-1}\varepsilon_t s_t^\top\,R(\tau),
\qquad
\widehat G_{\ell}
=\sum_{t=0}^{T-1}\Big(\Sigma^{-1}(\varepsilon_t\odot \varepsilon_t)-\mathbf 1\Big)\,R(\tau),
\]
and the stacked noise $\bar\varepsilon:=(\varepsilon_0^\top,\dots,\varepsilon_{T-1}^\top)^\top\sim\mathcal N(0,\overline\Sigma)$
with $\overline\Sigma:=I_T\otimes\Sigma$.

% =========================================================
% Expand both second moments
% =========================================================

\paragraph{Second moment expansions.}
Using $\|AB\|_{\mathrm{F}}^2=\tr(B^\top A^\top A B)$,
\[
\|\widehat G_K\|_{\mathrm{F}}^2
=\Big\|\sum_{t=0}^{T-1}\Sigma^{-1}\varepsilon_t s_t^\top\Big\|_{\mathrm{F}}^2\,R(\tau)^2
=\sum_{t,p=0}^{T-1}\big(\varepsilon_t^\top\Sigma^{-2}\varepsilon_p\big)\,(s_t^\top s_p)\,R(\tau)^2.
\]
For the log-std part, define
\[
g_t := \Sigma^{-1}(\varepsilon_t\odot \varepsilon_t)-\mathbf 1\in\mathbb R^m.
\]
Then
\[
\|\widehat G_{\ell}\|_2^2
=\Big\|\sum_{t=0}^{T-1} g_t\Big\|_2^2 R(\tau)^2
=\sum_{t,p=0}^{T-1}(g_t^\top g_p)\,R(\tau)^2.
\]

Using the lifted decomposition $R(\tau)=-(x+2y+z)$ with
\[
x=s_0^\top\overline M_{ss}s_0,\qquad
y=s_0^\top\overline M_{se}\bar\varepsilon,\qquad
z=\bar\varepsilon^\top\overline M_{ee}\bar\varepsilon,
\]
we have
\begin{equation}\label{eq:R2_expand_full}
R(\tau)^2=(x+2y+z)^2 = x^2+z^2+2xz+4y^2 + 4xy + 4yz .
\end{equation}
Therefore
\begin{equation}\label{eq:EGnorm_expand_tp}
\mathbb E\big[\|\widehat G_K\|_{\mathrm{F}}^2\big]
=\sum_{t,p=0}^{T-1}\mathbb E\Big[(\varepsilon_t^\top\Sigma^{-2}\varepsilon_p)(s_t^\top s_p)\,
\big(x^2+z^2+2xz+4y^2+4xy+4yz\big)\Big].
\end{equation}

\paragraph{Even-only reduction for the $\ell$-part.}
Note that $g_t^\top g_p$ is an \emph{even} polynomial in $\bar\varepsilon$, while $xy$ and $yz$ have \emph{odd} total
degree in $\bar\varepsilon$; hence
\[
\mathbb E_{\bar\varepsilon} \big[(g_t^\top g_p)\,(4xy+4yz)\mid s_0\big]=0,
\]
and therefore
\begin{equation}\label{eq:EGell_even_only}
\mathbb E\big[\|\widehat G_{\ell}\|_2^2\big]
=\sum_{t,p=0}^{T-1}\mathbb E\Big[(g_t^\top g_p)\,\big(x^2+z^2+2xz+4y^2\big)\Big].
\end{equation}

% =========================================================
% Multi step gradient
% =========================================================

\paragraph{Multi-step gradient.}
Consider the original dynamics~\eqref{eq:lqr-dynamics-reward} (not the lifted one). Since $\mathbb E[s_0]=0$ and $\mathbb E[\varepsilon_t]=0$, it follows by induction that $\mathbb E[s_t]=0$ for all $t$. Let $P_t:=\text{Cov}(s_t)$. Using independence of $s_t$ and $\varepsilon_t$,
\[
P_{t+1}=\text{Cov}(Fs_t+B\varepsilon_t)=F P_t F^\top + B\Sigma B^\top,\qquad P_0=\Sigma_0.
\]

Therefore
\[
\mathbb E\big[s_{t+1}^\top Q_s\, s_{t+1}\big]=\tr(Q_s\, P_{t+1}).
\]
Also $a_t=K s_t+\varepsilon_t$ is zero-mean with
\[
\text{Cov}(a_t)=K P_t K^\top + \Sigma
\quad\Rightarrow\quad
\mathbb E\big[a_t^\top Q_a\, a_t\big]=\tr \big(Q_a(K P_t K^\top+\Sigma)\big)
= \tr(Q_a K P_t K^\top)+\tr(Q_a\Sigma).
\]
Taking expectations in $r(s_{t+1},a_t)=-(s_{t+1}^\top Q_s\, s_{t+1}+a_t^\top Q_a\, a_t)$ and summing over $t=0,\dots,T-1$ with discount $\gamma^t$ yields
\begin{align}\label{eq:lqr-objective-matrix-form}
  J_T(K,\Sigma)
= -\sum_{t=0}^{T-1}\gamma^{t} \left(\tr(Q_s\,P_{t+1})
+\tr(Q_a\,K P_t K^\top)+\tr(Q_a\Sigma)\right).
\end{align}
To take the gradient of \eqref{eq:lqr-objective-matrix-form}, we use a Lagrangian to enforce the covariance dynamics $P_{t+1}=F P_t F^\top + B\Sigma B^\top$.
Introduce multipliers $\{\Lambda_{t+1}\}_{t=0}^{T-1}$ and set
\[
\mathcal L
= -\sum_{t=0}^{T-1}\gamma^{t} \left(\tr(Q_s\,P_{t+1})+\tr(Q_a\,K P_t K^\top)+\tr(Q_a\Sigma)\right)
\;+\;\sum_{t=0}^{T-1}\gamma^{t+1}\tr \Big(\Lambda_{t+1}\,(P_{t+1}-F P_t F^\top - B\Sigma B^\top)\Big).
\]

(i) Stationarity \wrt  $P_t$.
For $t=1,\dots,T-1$, the matrix $P_t$ appears in the terms at indices $t-1$ and $t$:
\[
\partial_{P_t}\mathcal L:\quad
-\gamma^{t-1}Q_s\;-\gamma^{t}K^\top Q_a K \;+\;\gamma^{t}\Lambda_t\;-\;\gamma^{t+1}F^\top \Lambda_{t+1}F \;=\;0.
\]
Dividing by $\gamma^{t}$ yields the recursion
$\Lambda_t=\tfrac{1}{\gamma}Q_s+K^\top Q_a K+\gamma F^\top\Lambda_{t+1}F$.
For $t=T$, only the $(T-1)$-indexed terms contain $P_T$, giving
$-\gamma^{T-1}Q_s+\gamma^{T}\Lambda_T=0$, hence $\Lambda_T=\tfrac{1}{\gamma}Q_s$.

(ii) Envelope / stationarity \wrt  $K$ and $\Sigma$.
Only two parts of $\mathcal L$ depend explicitly on $K$:
the quadratic term in $a_t$ and the dynamics through $F=A+BK$.
Using $dF=B\,dK$, symmetricity of $P_t, \Lambda_t$, and standard identities for matrix differentials,
\[
d\,\tr(Q_a\,K P_t K^\top) \;=\; 2\,\tr \big((Q_a K P_t)^\top dK\big),
\quad
d\,\tr(\Lambda_{t+1} F P_t F^\top)
=\;2\,\tr \big((B^\top\Lambda_{t+1}F P_t)^\top dK\big).
\]
Therefore,
\[
\partial_K\mathcal L
= -\,2\sum_{t=0}^{T-1}\gamma^t\,(Q_a K P_t)
\;-\;2\sum_{t=0}^{T-1}\gamma^{t+1}\,(B^\top\Lambda_{t+1}F P_t)
= -\,2\sum_{t=0}^{T-1}\gamma^{t}\Big( Q_a K P_t \;+\; \gamma\, B^\top \Lambda_{t+1} F\, P_t \Big)
\]

Only two places depend on $\Sigma$:
the immediate term $-\gamma^t\tr(Q_a\Sigma)$ and the constraint term
$-\gamma^{t+1}\tr(\Lambda_{t+1}\,B\Sigma B^\top)$. Then
\[
\frac{\partial \mathcal L}{\partial \Sigma}
= -\sum_{t=0}^{T-1}\gamma^t Q_a \;-\; \sum_{t=0}^{T-1}\gamma^{t+1}\,B^\top\Lambda_{t+1}B = -\sum_{t=0}^{T-1}\gamma^t\Big( Q_a \;+\; \gamma\, B^\top \Lambda_{t+1} B \Big),
\]
Therefore, the gradient with respect to the log-std parameter \(\ell\) is
\begin{equation}
\nabla_\ell J_T
\;=\;
-2\,\diag{\Sigma\sum_{t=0}^{T-1}\gamma^t\Big( Q_a \;+\; \gamma\, B^\top \Lambda_{t+1} B \Big)}
\;\;\in\mathbb{R}^{m}.
\end{equation}

% =========================================================
% K PART
% =========================================================

\paragraph{Variance of $K$-part: selectors and lifted state expansion.}
Introduce block selectors $\Pi_t$ such that $\varepsilon_t=\Pi_t\bar\varepsilon$. Then
\[
\varepsilon_t^\top\Sigma^{-2}\varepsilon_p
=\bar\varepsilon^\top W^K_{tp}\bar\varepsilon,
\qquad
W^K_{tp}:=\Pi_t^\top\Sigma^{-2}\Pi_p.
\]
Also, with $s_t=S_t\bar s$ and $\bar s=\mathcal F_S s_0+\mathcal F_E\bar\varepsilon$, define
\[
F_t:=S_t\mathcal F_S,\qquad E_t:=S_t\mathcal F_E,
\]
so that $s_t=F_t s_0+E_t\bar\varepsilon$ and hence
\begin{equation}\label{eq:stsp_full_expand}
s_t^\top s_p
=s_0^\top(F_t^\top F_p)s_0
+s_0^\top(F_t^\top E_p+F_p^\top E_t)\bar\varepsilon
+\bar\varepsilon^\top(E_t^\top E_p)\bar\varepsilon.
\end{equation}
Define
\[
G_{tp}:=F_t^\top F_p,\qquad
J_{tp}:=F_t^\top E_p+F_p^\top E_t,\qquad
U_{tp}:=E_t^\top E_p.
\]

\paragraph{Parity bookkeeping after conditioning on $s_0$ ($K$-part).}
Condition on $s_0$. Since $s_0 \perp\!\!\!\perp \bar\varepsilon$ and $\bar\varepsilon$ is centered Gaussian,
only even total degree polynomials in $\bar\varepsilon$ survive inside $\mathbb E_{\bar\varepsilon}[\cdot\mid s_0]$.
Note
\[
\bar\varepsilon^\top W^K_{tp}\bar\varepsilon \text{ is degree }2,\quad
s_0^\top G_{tp}s_0 \text{ is degree }0,\quad
s_0^\top J_{tp}\bar\varepsilon \text{ is degree }1,\quad
\bar\varepsilon^\top U_{tp}\bar\varepsilon \text{ is degree }2,
\]
and among \eqref{eq:R2_expand_full},
\[
x^2 \text{ degree }0,\quad
z^2 \text{ degree }4,\quad
xz \text{ degree }2,\quad
y^2 \text{ degree }2,\quad
xy \text{ degree }1,\quad
yz \text{ degree }3.
\]
Therefore, for the \emph{even} part $x^2+z^2+2xz+4y^2$, the linear term
$s_0^\top J_{tp}\bar\varepsilon$ in \eqref{eq:stsp_full_expand} yields an odd integrand and vanishes, so
\[
s_t^\top s_p \;\to\; s_0^\top G_{tp}s_0+\bar\varepsilon^\top U_{tp}\bar\varepsilon
\qquad\text{when multiplying}\qquad (x^2+z^2+2xz+4y^2).
\]
For the \emph{odd} part $4xy+4yz$, only the linear term survives (odd$\times$odd becomes even), hence
\[
s_t^\top s_p \;\to\; s_0^\top J_{tp}\bar\varepsilon
\qquad\text{when multiplying}\qquad (4xy+4yz).
\]

\paragraph{Six-term decomposition ($K$-part).}
Substituting $\varepsilon_t^\top\Sigma^{-2}\varepsilon_p=\bar\varepsilon^\top W^K_{tp}\bar\varepsilon$ and using the parity reductions,
\eqref{eq:EGnorm_expand_tp} becomes
\[
\mathbb E\big[\|\widehat G_K\|_{\mathrm{F}}^2\big]
=\sum_{t,p=0}^{T-1}\big(E^{K}_{1,tp}+E^{K}_{2,tp}+E^{K}_{3,tp}+E^{K}_{4,tp}+E^{K}_{5,tp}+E^{K}_{6,tp}\big),
\]
where the \emph{even-part} contributions are
\begin{align*}
E^{K}_{1,tp}&:=\mathbb E \left[(\bar\varepsilon^\top W^K_{tp}\bar\varepsilon)\Big(s_0^\top G_{tp}s_0+\bar\varepsilon^\top U_{tp}\bar\varepsilon\Big)\,x^2\right],\\
E^{K}_{2,tp}&:=\mathbb E \left[(\bar\varepsilon^\top W^K_{tp}\bar\varepsilon)\Big(s_0^\top G_{tp}s_0+\bar\varepsilon^\top U_{tp}\bar\varepsilon\Big)\,z^2\right],\\
E^{K}_{3,tp}&:=2\,\mathbb E \left[(\bar\varepsilon^\top W^K_{tp}\bar\varepsilon)\Big(s_0^\top G_{tp}s_0+\bar\varepsilon^\top U_{tp}\bar\varepsilon\Big)\,x z\right],\\
E^{K}_{4,tp}&:=4\,\mathbb E \left[(\bar\varepsilon^\top W^K_{tp}\bar\varepsilon)\Big(s_0^\top G_{tp}s_0+\bar\varepsilon^\top U_{tp}\bar\varepsilon\Big)\,y^2\right],
\end{align*}
and the \emph{odd-part} contributions are
\begin{equation}\label{eq:EtpFtp_def}
E^{K}_{5,tp}:=4\,\mathbb E \left[(\bar\varepsilon^\top W^K_{tp}\bar\varepsilon)\,(s_0^\top J_{tp}\bar\varepsilon)\,x\,y\right],\qquad
E^{K}_{6,tp}:=4\,\mathbb E \left[(\bar\varepsilon^\top W^K_{tp}\bar\varepsilon)\,(s_0^\top J_{tp}\bar\varepsilon)\,y\,z\right].
\end{equation}

\noindent\textbf{Term $E^{K}_{1,tp}$ ($x^2$-part).}
Since $x$ depends only on $s_0$, we can factor expectations:
\begin{align}\label{eq:E1tp_final}
E^{K}_{1,tp}
&=\mathbb E \left[(\bar\varepsilon^\top W^K_{tp}\bar\varepsilon)(s_0^\top G_{tp}s_0)\,x^2\right]
+\mathbb E \left[(\bar\varepsilon^\top W^K_{tp}\bar\varepsilon)(\bar\varepsilon^\top U_{tp}\bar\varepsilon)\,x^2\right] \nonumber \\
&=\mathrm{IS}_{\Sigma_0}\big(\symfun(G_{tp}),\symfun(\overline M_{ss}),\symfun(\overline M_{ss})\big)\;
  \mathrm{IS}_{\overline\Sigma}\big(\symfun(W^K_{tp})\big) \nonumber \\
&\quad+
  \mathrm{IS}_{\Sigma_0} \big(\symfun(\overline M_{ss}),\symfun(\overline M_{ss})\big)\;
  \mathrm{IS}_{\overline\Sigma}\big(\symfun(W^K_{tp}),\symfun(U_{tp})\big).
\end{align}

\medskip
\noindent\textbf{Term $E^{K}_{2,tp}$ ($z^2$-part).}
\begin{align}\label{eq:E2tp_final}
E^{K}_{2,tp}
&=\mathrm{IS}_{\Sigma_0}\big(\symfun(G_{tp})\big)\;
  \mathrm{IS}_{\overline\Sigma}\big(\symfun(W^K_{tp}),\symfun(\overline M_{ee}),\symfun(\overline M_{ee})\big) \nonumber \\
&\quad+
  \mathrm{IS}_{\overline\Sigma}\big(\symfun(W^K_{tp}),\symfun(U_{tp}),\symfun(\overline M_{ee}),\symfun(\overline M_{ee})\big).
\end{align}

\medskip
\noindent\textbf{Term $E^{K}_{3,tp}$ ($2xz$-part).}
\begin{align}\label{eq:E3tp_final}
E^{K}_{3,tp}
&=2\,\mathrm{IS}_{\Sigma_0}\big(\symfun(G_{tp}),\symfun(\overline M_{ss})\big)\;
     \mathrm{IS}_{\overline\Sigma}\big(\symfun(W^K_{tp}),\symfun(\overline M_{ee})\big) \nonumber \\
&\quad+
  2\,\mathrm{IS}_{\Sigma_0}\big(\symfun(\overline M_{ss})\big)\;
     \mathrm{IS}_{\overline\Sigma}\big(\symfun(W^K_{tp}),\symfun(U_{tp}),\symfun(\overline M_{ee})\big).
\end{align}

\medskip
\noindent\textbf{Term $E^{K}_{4,tp}$ ($4y^2$-part).}
Since $y=s_0^\top\overline M_{se}\bar\varepsilon$, we have
\[
y^2=\bar\varepsilon^\top H(s_0)\bar\varepsilon,\qquad
H(s_0):=\overline M_{se}^\top s_0 s_0^\top\overline M_{se}.
\]
Then
\begin{align}
E^{K}_{4,tp}
&=4\,\mathbb E \left[(\bar\varepsilon^\top W^K_{tp}\bar\varepsilon)\Big(s_0^\top G_{tp}s_0+\bar\varepsilon^\top U_{tp}\bar\varepsilon\Big)\,(\bar\varepsilon^\top H(s_0)\bar\varepsilon)\right]\nonumber\\
&=4\,\mathbb E_{s_0} \Bigg[\;
(s_0^\top G_{tp}s_0)\underbrace{\mathbb E_{\bar\varepsilon}\left[(\bar\varepsilon^\top W^K_{tp}\bar\varepsilon)(\bar\varepsilon^\top H(s_0)\bar\varepsilon)\,\middle|\,s_0\right]}_{=:T_1(s_0)}
+\underbrace{\mathbb E_{\bar\varepsilon} \left[(\bar\varepsilon^\top W^K_{tp}\bar\varepsilon)(\bar\varepsilon^\top U_{tp}\bar\varepsilon)(\bar\varepsilon^\top H(s_0)\bar\varepsilon)\,\middle|\,s_0\right]}_{=:T_2(s_0)}
\;\Bigg].\label{eq:Dtp_cond_expand}
\end{align}
By Lemma~\ref{lem:IS_shorthand} (the $k=2$ and $k=3$ cases) with $\Omega=\overline\Sigma$,
\begin{align}
T_1(s_0)
&=\mathrm{IS}_{\overline\Sigma}\big(\symfun(W^K_{tp}),H(s_0)\big)\notag\\
&=\tr(\overline\Sigma\,\symfun(W^K_{tp}))\tr(\overline\Sigma\,H(s_0))
+2\,\tr(\overline\Sigma\,\symfun(W^K_{tp})\,\overline\Sigma\,H(s_0)),\label{eq:T1_expand}\\
T_2(s_0)
&=\mathrm{IS}_{\overline\Sigma}\big(\symfun(W^K_{tp}),\symfun(U_{tp}),H(s_0)\big)\notag\\
&=\tr(\overline\Sigma\,\symfun(W^K_{tp}))\tr(\overline\Sigma\,\symfun(U_{tp}))\tr(\overline\Sigma\,H(s_0))
+2\,\tr(\overline\Sigma\,\symfun(W^K_{tp})\,\overline\Sigma\,\symfun(U_{tp}))\tr(\overline\Sigma\,H(s_0))\notag\\
&\quad
+2\,\tr(\overline\Sigma\,\symfun(W^K_{tp})\,\overline\Sigma\,H(s_0))\tr(\overline\Sigma\,\symfun(U_{tp}))
+2\,\tr(\overline\Sigma\,\symfun(U_{tp})\,\overline\Sigma\,H(s_0))\tr(\overline\Sigma\,\symfun(W^K_{tp}))\notag\\
&\quad
+8\,\tr(\overline\Sigma\,\symfun(W^K_{tp})\,\overline\Sigma\,\symfun(U_{tp})\,\overline\Sigma\,H(s_0)).\label{eq:T2_expand}
\end{align}
Introduce $\alpha_{tp}:=\tr(\overline\Sigma\,\symfun(W^K_{tp}))$, $\beta_{tp}:=\tr(\overline\Sigma\,\symfun(U_{tp}))$,
$\delta_{tp}:=\tr(\overline\Sigma\,\symfun(W^K_{tp})\,\overline\Sigma\,\symfun(U_{tp}))$, and the matrices
\[
V_0:=\overline M_{se}\,\overline\Sigma\,\overline M_{se}^\top,\qquad
V_{tp}:=\overline M_{se}\,\overline\Sigma\,\symfun(W^K_{tp})\,\overline\Sigma\,\overline M_{se}^\top,
\]
\[
\widetilde V_{tp}:=\overline M_{se}\,\overline\Sigma\,\symfun(U_{tp})\,\overline\Sigma\,\overline M_{se}^\top,\qquad
\widehat V_{tp}:=\overline M_{se}\,\overline\Sigma\,\symfun(W^K_{tp})\,\overline\Sigma\,\symfun(U_{tp})\,\overline\Sigma\,\overline M_{se}^\top.
\]
Using $H(s_0)=\overline M_{se}^\top s_0 s_0^\top\overline M_{se}$ and cyclicity of trace, we obtain
\[
\tr(\overline\Sigma\,H(s_0))=s_0^\top V_0 s_0,\qquad
\tr(\overline\Sigma\,\symfun(W^K_{tp})\,\overline\Sigma\,H(s_0))=s_0^\top V_{tp} s_0,
\]
\[
\tr(\overline\Sigma\,\symfun(U_{tp})\,\overline\Sigma\,H(s_0))=s_0^\top \widetilde V_{tp} s_0,\qquad
\tr(\overline\Sigma\,\symfun(W^K_{tp})\,\overline\Sigma\,\symfun(U_{tp})\,\overline\Sigma\,H(s_0))=s_0^\top \widehat V_{tp} s_0.
\]
Substituting into \eqref{eq:T1_expand}--\eqref{eq:T2_expand} yields
\[
T_1(s_0)=\alpha_{tp}(s_0^\top V_0 s_0)+2(s_0^\top V_{tp}s_0),
\]
\[
T_2(s_0)
=(\alpha_{tp}\beta_{tp}+2\delta_{tp})(s_0^\top V_0 s_0)
+2\beta_{tp}(s_0^\top V_{tp}s_0)
+2\alpha_{tp}(s_0^\top \widetilde V_{tp}s_0)
+8(s_0^\top \widehat V_{tp}s_0).
\]
Taking $\mathbb E_{s_0}$ term-by-term and using Lemma~\ref{lem:IS_shorthand} with $\Omega=\Sigma_0$ gives
\begin{align}\label{eq:Dtp_final}
E^{K}_{4,tp}
&=
4\Big(
\alpha_{tp}\,\mathrm{IS}_{\Sigma_0}(\symfun(G_{tp}),V_0)
+2\,\mathrm{IS}_{\Sigma_0}(\symfun(G_{tp}),V_{tp})
+(\alpha_{tp}\beta_{tp}+2\delta_{tp})\,\mathrm{IS}_{\Sigma_0}(V_0) \nonumber\\
&+2\beta_{tp}\,\mathrm{IS}_{\Sigma_0}(V_{tp})
+2\alpha_{tp}\,\mathrm{IS}_{\Sigma_0}(\widetilde V_{tp})
+8\,\mathrm{IS}_{\Sigma_0} \big(\symfun(\widehat V_{tp})\big)
\Big).
\end{align}

\medskip
\noindent\textbf{Term $E^{K}_{5,tp}$ ($4xy$-part).}
Recall $y=s_0^\top\overline M_{se}\bar\varepsilon$. Define
\[
L_{tp}(s_0):=J_{tp}^\top s_0 s_0^\top \overline M_{se}\in\mathbb R^{mT\times mT}.
\]
Then
\[
(s_0^\top J_{tp}\bar\varepsilon)\,(s_0^\top \overline M_{se}\bar\varepsilon)
=\bar\varepsilon^\top L_{tp}(s_0)\bar\varepsilon,
\]
and since $x=s_0^\top\overline M_{ss}s_0$ depends only on $s_0$, conditioning on $s_0$ gives
\begin{align}
E^{K}_{5,tp}
&=4\,\mathbb E_{s_0} \left[
x\;\mathbb E_{\bar\varepsilon} \left[(\bar\varepsilon^\top W^K_{tp}\bar\varepsilon)\,(\bar\varepsilon^\top L_{tp}(s_0)\bar\varepsilon)\,\middle|\,s_0\right]\right]\nonumber\\
&=4\,\mathbb E_{s_0} \left[x\;\mathrm{IS}_{\overline\Sigma} \big(\symfun(W^K_{tp}),\symfun(L_{tp}(s_0))\big)\right].\label{eq:Etp_cond}
\end{align}
Introduce $\alpha_{tp}:=\tr(\overline\Sigma\,\symfun(W^K_{tp}))$ and define
\[
V^{(0)}_{tp}:=\symfun \Big(\overline M_{se}\,\overline\Sigma\,J_{tp}^\top\Big),\qquad
V^{(1)}_{tp}:=\symfun \Big(\overline M_{se}\,\overline\Sigma\,\symfun(W^K_{tp})\,\overline\Sigma\,J_{tp}^\top\Big).
\]
Then the same trace-to-quadratic-form yields
\begin{equation}\label{eq:Etp_final}
E^{K}_{5,tp}
=
4\Big(
\alpha_{tp}\,\mathrm{IS}_{\Sigma_0} \big(\symfun(\overline M_{ss}),V^{(0)}_{tp}\big)
+2\,\mathrm{IS}_{\Sigma_0} \big(\symfun(\overline M_{ss}),V^{(1)}_{tp}\big)
\Big).
\end{equation}

\medskip
\noindent\textbf{Term $E^{K}_{6,tp}$ ($4yz$-part).}
Using the same $L_{tp}(s_0)$ and $z=\bar\varepsilon^\top\overline M_{ee}\bar\varepsilon$, conditioning on $s_0$ gives
\begin{align}
E^{K}_{6,tp}
&=4\,\mathbb E_{s_0} \left[
\mathrm{IS}_{\overline\Sigma} \Big(\symfun(W^K_{tp}),\symfun(\overline M_{ee}),\symfun(L_{tp}(s_0))\Big)\right].
\label{eq:Ftp_cond}
\end{align}
Define
\[
\alpha_{tp}:=\tr(\overline\Sigma\,\symfun(W^K_{tp})),\qquad
\mu:=\tr(\overline\Sigma\,\symfun(\overline M_{ee})),\qquad
\kappa_{tp}:=\tr(\overline\Sigma\,\symfun(W^K_{tp})\,\overline\Sigma\,\symfun(\overline M_{ee})).
\]
With
\[
A^{(0)}:=\overline\Sigma,\quad
A^{(1)}:=\overline\Sigma\,\symfun(W^K_{tp})\,\overline\Sigma,\quad
A^{(2)}:=\overline\Sigma\,\symfun(\overline M_{ee})\,\overline\Sigma,\quad
A^{(3)}:=\symfun(\overline\Sigma\,\symfun(W^K_{tp})\,\overline\Sigma\,\symfun(\overline M_{ee})\,\overline\Sigma),
\]
and
\[
V^{(j)}_{tp}:=\symfun \Big(\overline M_{se}\,A^{(j)}\,J_{tp}^\top\Big)\quad (j=0,1,2,3),
\]
the trace-to-quadratic-form conversion gives
\begin{equation}\label{eq:Ftp_final}
E^{K}_{6,tp}
=
4\Big(
(\alpha_{tp}\mu+2\kappa_{tp})\,\mathrm{IS}_{\Sigma_0} \big(V^{(0)}_{tp}\big)
+2\mu\,\mathrm{IS}_{\Sigma_0} \big(V^{(1)}_{tp}\big)
+2\alpha_{tp}\,\mathrm{IS}_{\Sigma_0} \big(V^{(2)}_{tp}\big)
+8\,\mathrm{IS}_{\Sigma_0} \big(V^{(3)}_{tp}\big)
\Big).
\end{equation}

% =========================================================
% ell PART
% =========================================================

\paragraph{$\ell$-part: lifted representation of $g_t^\top g_p$.}
Introduce $\Pi_t$ such that $\varepsilon_t=\Pi_t\bar\varepsilon$.
Write $D:=\Sigma^{-1}=\mathrm{diag}(d_1,\dots,d_m)$.
Then
\[
g_t^\top g_p
=\big(D(\varepsilon_t\odot\varepsilon_t)-\mathbf 1\big)^\top\big(D(\varepsilon_p\odot\varepsilon_p)-\mathbf 1\big)
=(\varepsilon_t\odot\varepsilon_t)^\top D^2(\varepsilon_p\odot\varepsilon_p)
-\varepsilon_t^\top D\varepsilon_t-\varepsilon_p^\top D\varepsilon_p+m.
\]
To express the quartic term via quadratic forms, define for each $i\in\{1,\dots,m\}$
\[
Q_{t,i}:=\Pi_t^\top e_ie_i^\top \Pi_t\in\mathbb R^{mT\times mT}
\quad\Rightarrow\quad
\varepsilon_{t,i}^2=\bar\varepsilon^\top Q_{t,i}\bar\varepsilon,
\]
and define
\[
W^\ell_t:=\Pi_t^\top D \Pi_t\in\mathbb R^{mT\times mT}
\quad\Rightarrow\quad
\varepsilon_t^\top D\varepsilon_t=\bar\varepsilon^\top W^\ell_t\bar\varepsilon.
\]
Then
\begin{equation}\label{eq:gtgp_lifted}
g_t^\top g_p
=\sum_{i=1}^m d_i^2\,
(\bar\varepsilon^\top Q_{t,i}\bar\varepsilon)\,(\bar\varepsilon^\top Q_{p,i}\bar\varepsilon)
-(\bar\varepsilon^\top W^\ell_t\bar\varepsilon)-(\bar\varepsilon^\top W^\ell_p\bar\varepsilon)+m.
\end{equation}

\paragraph{Four-term decomposition ($\ell$-part).}
From \eqref{eq:EGell_even_only},
\[
\mathbb E\big[\|\widehat G_{\ell}\|_2^2\big]
=\sum_{t,p=0}^{T-1}\big(E^{\ell}_{1,tp}+E^{\ell}_{2,tp}+E^{\ell}_{3,tp}+E^{\ell}_{4,tp}\big),
\]
where
\[
E^{\ell}_{1,tp}:=\mathbb E[(g_t^\top g_p)x^2],\quad
E^{\ell}_{2,tp}:=\mathbb E[(g_t^\top g_p)z^2],\quad
E^{\ell}_{3,tp}:=2\,\mathbb E[(g_t^\top g_p)xz],\quad
E^{\ell}_{4,tp}:=4\,\mathbb E[(g_t^\top g_p)y^2].
\]

\medskip
\noindent\textbf{Term $E^{\ell}_{1,tp}$ ($x^2$-part).}
Since $x$ depends only on $s_0$ and $g_t^\top g_p$ depends only on $\bar\varepsilon$,
\[
E^{\ell}_{1,tp}=\mathbb E[x^2]\;\mathbb E[g_t^\top g_p]
=\mathrm{IS}_{\Sigma_0} \big(\symfun(\overline M_{ss}),\symfun(\overline M_{ss})\big)\;\mathbb E[g_t^\top g_p].
\]
Under $\overline\Sigma=I_T\otimes\Sigma$, $\{g_t\}$ are independent across time and $\mathbb E[g_t]=0$, hence
\[
\mathbb E[g_t^\top g_p]=\mathbb E[\|g_t\|_2^2]\mathbf 1_{\{t=p\}}=2m\,\mathbf 1_{\{t=p\}},
\]
so
\begin{equation}\label{eq:Atp_ell_final}
E^{\ell}_{1,tp}
=2m\,\mathbf 1_{\{t=p\}}\;\mathrm{IS}_{\Sigma_0} \big(\symfun(\overline M_{ss}),\symfun(\overline M_{ss})\big).
\end{equation}

\medskip
\noindent\textbf{Term $E^{\ell}_{2,tp}$ ($z^2$-part).}
Using \eqref{eq:gtgp_lifted} and $z=\bar\varepsilon^\top\overline M_{ee}\bar\varepsilon$,
\begin{align}
E^{\ell}_{2,tp}
&=\sum_{i=1}^m d_i^2\,
\mathrm{IS}_{\overline\Sigma} \big(\symfun(Q_{t,i}),\symfun(Q_{p,i}),\symfun(\overline M_{ee}),\symfun(\overline M_{ee})\big)\notag\\
&\quad-\mathrm{IS}_{\overline\Sigma} \big(\symfun(W^\ell_t),\symfun(\overline M_{ee}),\symfun(\overline M_{ee})\big)
-\mathrm{IS}_{\overline\Sigma} \big(\symfun(W^\ell_p),\symfun(\overline M_{ee}),\symfun(\overline M_{ee})\big)\notag\\
&\quad+m\,\mathrm{IS}_{\overline\Sigma} \big(\symfun(\overline M_{ee}),\symfun(\overline M_{ee})\big).
\label{eq:Btp_ell_final}
\end{align}

\medskip
\noindent\textbf{Term $E^{\ell}_{3,tp}$ ($2xz$-part).}
Since $x$ depends only on $s_0$,
\[
E^{\ell}_{3,tp}=2\,\mathbb E[x]\;\mathbb E[(g_t^\top g_p)z]
=2\,\mathrm{IS}_{\Sigma_0} \big(\symfun(\overline M_{ss})\big)\;\mathbb E[(g_t^\top g_p)z].
\]
By \eqref{eq:gtgp_lifted},
\begin{align}
\mathbb E[(g_t^\top g_p)z]
&=\sum_{i=1}^m d_i^2\,
\mathrm{IS}_{\overline\Sigma} \big(\symfun(Q_{t,i}),\symfun(Q_{p,i}),\symfun(\overline M_{ee})\big)\notag\\
&\quad-\mathrm{IS}_{\overline\Sigma} \big(\symfun(W^\ell_t),\symfun(\overline M_{ee})\big)
-\mathrm{IS}_{\overline\Sigma} \big(\symfun(W^\ell_p),\symfun(\overline M_{ee})\big)
+m\,\mathrm{IS}_{\overline\Sigma} \big(\symfun(\overline M_{ee})\big),
\end{align}
hence
\begin{align}\label{eq:Ctp_ell_final}
E^{\ell}_{3,tp}
&=2\,\mathrm{IS}_{\Sigma_0} \big(\symfun(\overline M_{ss})\big)\Big[
\sum_{i=1}^m d_i^2\,
\mathrm{IS}_{\overline\Sigma} \big(\symfun(Q_{t,i}),\symfun(Q_{p,i}),\symfun(\overline M_{ee})\big)
-\mathrm{IS}_{\overline\Sigma} \big(\symfun(W^\ell_t),\symfun(\overline M_{ee})\big)\notag\\
&\hspace{3.0cm}
-\mathrm{IS}_{\overline\Sigma} \big(\symfun(W^\ell_p),\symfun(\overline M_{ee})\big)
+m\,\mathrm{IS}_{\overline\Sigma} \big(\symfun(\overline M_{ee})\big)\Big].
\end{align}

\medskip
\noindent\textbf{Term $E^{\ell}_{4,tp}$ ($4y^2$-part).}
Recall $y=s_0^\top\overline M_{se}\bar\varepsilon$ and define
\[
H(s_0):=\overline M_{se}^\top s_0 s_0^\top\overline M_{se}
\quad\Rightarrow\quad
y^2=\bar\varepsilon^\top H(s_0)\bar\varepsilon .
\]
Conditioning on $s_0$ and using \eqref{eq:gtgp_lifted} yields
\begin{align}
E^{\ell}_{4,tp}
&=4\,\mathbb E_{s_0}\Big[
\sum_{i=1}^m d_i^2\,\mathrm{IS}_{\overline\Sigma} \big(\symfun(Q_{t,i}),\symfun(Q_{p,i}),H(s_0)\big)
-\mathrm{IS}_{\overline\Sigma} \big(\symfun(W^\ell_t),H(s_0)\big)\notag\\
&\hspace{3.6cm}
-\mathrm{IS}_{\overline\Sigma} \big(\symfun(W^\ell_p),H(s_0)\big)
+m\,\mathrm{IS}_{\overline\Sigma} \big(H(s_0)\big)\Big].
\label{eq:Dtp_ell_cond}
\end{align}

Introduce $c(s_0):=\tr(\overline\Sigma\,H(s_0))$. Using cyclicity of trace,
\begin{equation}\label{eq:cs0_V0}
c(s_0)=s_0^\top V_0 s_0,\qquad V_0:=\overline M_{se}\,\overline\Sigma\,\overline M_{se}^\top.
\end{equation}
Define also
\[
V_{t,i}:=\overline M_{se}\,\overline\Sigma\,\symfun(Q_{t,i})\,\overline\Sigma\,\overline M_{se}^\top,\qquad
V_{tp,i}:=\symfun \Big(\overline M_{se}\,\overline\Sigma\,\symfun(Q_{t,i})\,\overline\Sigma\,\symfun(Q_{p,i})\,\overline\Sigma\,\overline M_{se}^\top\Big),
\]
so that
\begin{equation}\label{eq:trace_to_quad_forms}
\tr(\overline\Sigma\,\symfun(Q_{t,i})\,\overline\Sigma\,H(s_0))=s_0^\top V_{t,i}s_0,\qquad
\tr(\overline\Sigma\,\symfun(Q_{t,i})\,\overline\Sigma\,\symfun(Q_{p,i})\,\overline\Sigma\,H(s_0))=s_0^\top V_{tp,i}s_0.
\end{equation}

\emph{Step 1: expand the $k=3$ term.}
By Lemma~\ref{lem:IS_shorthand} (the $k=3$ case) with $\Omega=\overline\Sigma$,
\begin{align}
\mathrm{IS}_{\overline\Sigma} \big(\symfun(Q_{t,i}),\symfun(Q_{p,i}),H(s_0)\big)
&=\tr(\overline\Sigma\,Q_{t,i})\,\tr(\overline\Sigma\,Q_{p,i})\,c(s_0)
+2\,\tr(\overline\Sigma\,Q_{t,i}\,\overline\Sigma\,Q_{p,i})\,c(s_0)\notag\\
&\quad
+2\,\tr(\overline\Sigma\,Q_{t,i}\,\overline\Sigma\,H(s_0))\,\tr(\overline\Sigma\,Q_{p,i})
+2\,\tr(\overline\Sigma\,Q_{p,i}\,\overline\Sigma\,H(s_0))\,\tr(\overline\Sigma\,Q_{t,i})\notag\\
&\quad
+8\,\tr(\overline\Sigma\,Q_{t,i}\,\overline\Sigma\,Q_{p,i}\,\overline\Sigma\,H(s_0)).
\label{eq:IS3_QQH_expand}
\end{align}
Under $\overline\Sigma=I_T\otimes\Sigma$,
\[
\tr(\overline\Sigma\,Q_{t,i})=\sigma_i^2,\qquad
\tr(\overline\Sigma\,Q_{t,i}\,\overline\Sigma\,Q_{p,i})=\mathbf 1_{\{t=p\}}\sigma_i^4.
\]
Multiplying by $d_i^2=\sigma_i^{-4}$ and using \eqref{eq:trace_to_quad_forms} gives
\begin{align}
d_i^2\,\mathrm{IS}_{\overline\Sigma} \big(Q_{t,i},Q_{p,i},H(s_0)\big)
&=c(s_0)+2\,\mathbf 1_{\{t=p\}}\,c(s_0)
+2d_i\,s_0^\top V_{t,i}s_0
+2d_i\,s_0^\top V_{p,i}s_0
+8d_i^2\,s_0^\top V_{tp,i}s_0.
\label{eq:di2_IS3_simplified}
\end{align}
Summing \eqref{eq:di2_IS3_simplified} over $i$ yields
\begin{align}
\sum_{i=1}^m d_i^2\,\mathrm{IS}_{\overline\Sigma} \big(\symfun(Q_{t,i}),\symfun(Q_{p,i}),H(s_0)\big)
&= m\,c(s_0)+2m\,\mathbf 1_{\{t=p\}}\,c(s_0)
+2\,s_0^\top\Big(\sum_{i=1}^m d_i V_{t,i}\Big)s_0
+2\,s_0^\top\Big(\sum_{i=1}^m d_i V_{p,i}\Big)s_0\notag\\
&\quad
+8\,s_0^\top\Big(\sum_{i=1}^m d_i^2 V_{tp,i}\Big)s_0.
\label{eq:sum_di2_IS3}
\end{align}

\emph{Step 2: expand the $k=2$ terms.}
By Lemma~\ref{lem:IS_shorthand} (the $k=2$ case),
\begin{equation}\label{eq:IS2_WH}
\mathrm{IS}_{\overline\Sigma} \big(W^\ell_t,H(s_0)\big)
=\tr(\overline\Sigma\,W^\ell_t)\,c(s_0)+2\,\tr(\overline\Sigma\,W^\ell_t\,\overline\Sigma\,H(s_0)).
\end{equation}
Under $\overline\Sigma=I_T\otimes\Sigma$ and $W^\ell_t=\Pi_t^\top\Sigma^{-1}\Pi_t$,
\[
\tr(\overline\Sigma\,W^\ell_t)=\tr(\Sigma\,\Sigma^{-1})=m,
\]
and
\begin{equation}\label{eq:trace_WtH_to_quad}
\tr(\overline\Sigma\,W^\ell_t\,\overline\Sigma\,H(s_0))
=s_0^\top\Big(\overline M_{se}\,\overline\Sigma\,W^\ell_t\,\overline\Sigma\,\overline M_{se}^\top\Big)s_0.
\end{equation}
Thus
\begin{equation}\label{eq:IS2_WH_simplified}
\mathrm{IS}_{\overline\Sigma} \big(W^\ell_t,H(s_0)\big)
=m\,c(s_0)+2\,s_0^\top\Big(\overline M_{se}\,\overline\Sigma\,W^\ell_t\,\overline\Sigma\,\overline M_{se}^\top\Big)s_0,
\end{equation}
and similarly for $p$.

\emph{Step 3: cancellation via $\sum_i d_i Q_{t,i}=W^\ell_t$.}
Note that
\[
\sum_{i=1}^m d_i\,Q_{t,i}
=\Pi_t^\top\Big(\sum_{i=1}^m d_i e_ie_i^\top\Big)\Pi_t
=\Pi_t^\top \Sigma^{-1} \Pi_t
=W^\ell_t.
\]
Therefore
\begin{equation}\label{eq:key_identity_sumdiVti}
\sum_{i=1}^m d_i\,V_{t,i}
=\overline M_{se}\,\overline\Sigma\,W^\ell_t\,\overline\Sigma\,\overline M_{se}^\top,
\qquad
\sum_{i=1}^m d_i\,V_{p,i}
=\overline M_{se}\,\overline\Sigma\,W^\ell_p\,\overline\Sigma\,\overline M_{se}^\top.
\end{equation}

\emph{Step 4: assemble \eqref{eq:Dtp_ell_cond}.}
Substituting \eqref{eq:sum_di2_IS3}, \eqref{eq:IS2_WH_simplified}, and $\mathrm{IS}_{\overline\Sigma}(H(s_0))=c(s_0)$,
and using the cancellation \eqref{eq:key_identity_sumdiVti}, the bracket in \eqref{eq:Dtp_ell_cond} reduces to
\[
2m\,\mathbf 1_{\{t=p\}}\,c(s_0)+8\,s_0^\top\Big(\sum_{i=1}^m d_i^2 V_{tp,i}\Big)s_0.
\]
Hence
\[
E^{\ell}_{4,tp}
=4\,\mathbb E_{s_0} \left[2m\,\mathbf 1_{\{t=p\}}\,c(s_0)+8\,s_0^\top\Big(\sum_{i=1}^m d_i^2 V_{tp,i}\Big)s_0\right].
\]
Using Lemma~\ref{lem:IS_shorthand} (the $k=1$ case) with $\Omega=\Sigma_0$ yields
\begin{equation}\label{eq:Dtp_ell_final}
E^{\ell}_{4,tp}
=
8m\,\mathbf 1_{\{t=p\}}\;\mathrm{IS}_{\Sigma_0}(V_0)
\;+\;
32\sum_{i=1}^m d_i^2\,\mathrm{IS}_{\Sigma_0} \big(V_{tp,i}\big).
\end{equation}
Moreover, under $\overline\Sigma=I_T\otimes\Sigma$ the matrices $Q_{t,i}$ and $Q_{p,i}$ occupy disjoint time blocks when
$t\neq p$, which implies $V_{tp,i}=0$ for $t\neq p$.

\paragraph{Conclusion.}
The $K$-part decomposes as
\[
\mathbb E\big[\|\widehat G_K\|_{\mathrm{F}}^2\big]
=\sum_{t,p=0}^{T-1}\big(E^{K}_{1,tp}+E^{K}_{2,tp}+E^{K}_{3,tp}+E^{K}_{4,tp}+E^{K}_{5,tp}+E^{K}_{6,tp}\big),
\]
with  $E^{K}_{i,tp}, i = 1,...,6$ given by \eqref{eq:E1tp_final}, \eqref{eq:E2tp_final}, \eqref{eq:E3tp_final}, \eqref{eq:Dtp_final}, \eqref{eq:Etp_final}, \eqref{eq:Ftp_final},
The $\ell$-part decomposes as
\[
\mathbb E\big[\|\widehat G_{\ell}\|_2^2\big]
=\sum_{t,p=0}^{T-1}\big(E^{\ell}_{1,tp}+E^{\ell}_{2,tp}+E^{\ell}_{3,tp}+E^{\ell}_{4,tp}\big),
\]
with $E^{\ell}_{1,tp}$, $E^{\ell}_{2,tp}$, $E^{\ell}_{3,tp}$, $E^{\ell}_{4,tp}$ given by
\eqref{eq:Atp_ell_final}, \eqref{eq:Btp_ell_final}, \eqref{eq:Ctp_ell_final}, \eqref{eq:Dtp_ell_final}.
This completes the proof.
\end{proof}


\section{Proof of Corollary~\ref{cor:lqg_stationary}}
\label{sec:app-proof-lqg-stationary}

\begin{proof}
From Theorem~\ref{thm:lqr-variance-T-compose}, the mean gradient with respect to $\ell$ is
\[
\nabla_{\ell} J_T
=
-2\,\mathrm{diag} \Big(\Sigma\sum_{t=0}^{T-1}\gamma^t\big( Q_a + \gamma B^\top \Lambda_{t+1} B \big)\Big).
\]
Define
\[
M := \sum_{t=0}^{T-1}\gamma^t\big( Q_a + \gamma B^\top \Lambda_{t+1} B \big).
\]
Then $\nabla_\ell J_T=0$ is equivalent to $\mathrm{diag}(\Sigma M)=0$.

As $\Sigma=\mathrm{diag}(\sigma^2)$ is diagonal, for each $i\in[m]$,
\[
0=(\Sigma M)_{ii}=\Sigma_{ii}M_{ii}.
\]
Since $Q_a\succ 0$ and $\Lambda_{t+1}\succeq 0$, we have
$Q_a+\gamma B^\top \Lambda_{t+1}B \succ 0$ for every $t$, hence $M\succ 0$ and in particular
$M_{ii}>0$ for all $i$. Therefore $\Sigma_{ii}=0$ for all $i$, i.e.\ $\Sigma=\mathbf 0$.

Conversely, if $\Sigma=\mathbf 0$, then $\nabla_\ell J_T=0$ follows immediately from the formula above.
\end{proof}
